\documentclass[10pt]{article}
\usepackage[T1]{fontenc}
\usepackage[utf8]{inputenc}
\usepackage[french]{babel}
\usepackage{amsmath,amssymb,mathrsfs,mathpazo}
\usepackage{geometry}
\geometry{a4paper,margin=2cm}
\usepackage{graphicx,array,tabularx,tabularray,colortbl,multirow,multicol,booktabs,enumitem,pifont,ulem,fancybox,diagbox,cancel}
\usepackage{pstricks,pstricks-add,pst-plot,pst-tree,pst-node,pst-text,pst-func,pst-math,pst-bspline,pst-3dplot,pst-eucl,pst-solides3d}
\let\pscalculate\undefined
\usepackage{pst-bezier}
\usepackage{tikz,pgfplots,tkz-tab}
\usetikzlibrary{arrows,arrows.meta,patterns,calc,positioning,trees,decorations.markings,shapes.geometric,shapes.misc,shapes.arrows,intersections,decorations.pathmorphing,decorations.pathreplacing,backgrounds,shadows,babel}
\pgfplotsset{compat=1.18}
\usepgfplotslibrary{fillbetween}
\usepackage[np]{numprint}
\usepackage{listings,makecell,pas-tableur,twoopt}
\let\barre\undefined
\usepackage[pstricks]{tablvar}
\definecolor{darkgreen}{RGB}{0,100,0}
\providecommand{\red}{\color{red}}
\providecommand{\blue}{\color{blue}}
\providecommand{\green}{\color{green}}
\providecommand{\black}{\color{black}}
\providecommand{\white}{\color{white}}
\providecommand{\gray}{\color{gray}}
\providecommand{\cyan}{\color{cyan}}
\providecommand{\magenta}{\color{magenta}}
\providecommand{\yellow}{\color{yellow}}
\providecommand{\violet}{\color{violet}}
\providecommand{\orange}{\color{orange}}
\providecommand{\euro}{\texteuro}
\providecommand{\EUR}{\texteuro}
\providecommand{\ieme}{\textsuperscript{e}}
\providecommand{\degres}{\ensuremath{^\circ}}
\providecommand{\degree}{\ensuremath{^\circ}}
\providecommand{\vv}[1]{\overrightarrow{#1}}
\providecommand{\Coord}[1]{#1}
\providecommand{\decosix}{}
\providecommand{\starredbullet}{\ensuremath{\bullet}}
\pagestyle{empty}
\setlength{\parindent}{0pt}
\usepackage{adjustbox}\DeclareUnicodeCharacter{2236}{:}\DeclareUnicodeCharacter{2190}{\ensuremath{\leftarrow}}
\providecommand{\R}{\mathbb{R}}
\providecommand{\N}{\mathbb{N}}
\providecommand{\Z}{\mathbb{Z}}
\providecommand{\Q}{\mathbb{Q}}
\providecommand{\D}{\mathbb{D}}
\providecommand{\C}{\mathbb{C}}
\providecommand{\e}{{\mathrm{e}}}
\providecommand{\vect}[1]{\overrightarrow{#1}}
\providecommand{\vectt}[1]{\overrightarrow{\mathrm{#1}}}
\providecommand{\vvt}[1]{\overrightarrow{\mathrm{#1}}}
\providecommand{\np}[2][]{{#2~#1}}
\providecommand{\vv}[1]{{\overrightarrow{#1}}}
\renewcommand{\d}{{\mathrm{d}}}
\providecommand{\degres}{{^\circ}}
\providecommand{\degre}{{^\circ}}
\providecommand{\blue}{\color{blue}}
\providecommand{\red}{\color{red}}
\providecommand{\enskip}{\quad}
\providecommand{\llbracket}{\lbrack\!\lbrack}
\providecommand{\rrbracket}{\rbrack\!\rbrack}
\providecommand{\euro}{{\text{€}}}
\providecommand{\no}{{\text{n°}}}
\providecommand{\interval}[2]{{[#1;#2]}}
\providecommand{\intervFO}[2]{{[#1;#2[}}
\providecommand{\dotfill}{\ldots}

\providecommand{\ds}{\displaystyle}
\providecommand{\textding}[1]{#1}
\providecommand{\Oij}{\(\mathrm{O};\vec{\imath},\vec{\jmath}\)}
\providecommand{\Oijk}{\(\mathrm{O};\vec{\imath},\vec{\jmath},\vec{k}\)}
\providecommand{\Ouv}{\(\mathrm{O};\vec{u},\vec{v}\)}
\providecommand{\intervFF}[2]{\ensuremath{[#1;#2]}}
\renewcommand{\euro}{\ifmmode\text{\texteuro}\else\texteuro\fi}
\DeclareUnicodeCharacter{00B0}{\ifmmode^{\circ}\else\textdegree\fi}
\DeclareUnicodeCharacter{20AC}{\ifmmode\text{\texteuro}\else\texteuro\fi}
% Énoncés : Éducation nationale. Transcription et corrigés : APMEP, Denis Vergès et contributeurs. Découpage et classement : CoopMaths, Sébastien Lozano, Éric Elter et contributeurs. Adaptation de la présentation : Maths Explorer.
% CC BY-SA : https://coopmaths.fr/www/about/
% https://coopmaths.fr/alea/static/bac/2021/tex/bac_2021_06_metropole_4_cor.tex
\begin{document}
\noindent\textit{APMEP / CoopMaths. CC BY-SA. Présentation adaptée par Maths Explorer.}\par\medskip
\vspace{0.75cm}
\begin{tabular}[]{|l|}
\hline
Principaux domaines abordés :\\

Géométrie de l'espace rapporté  à un repère orthonormé ; orthogonalité dans l'espace\\
\hline
\end{tabular}

\vspace{0.5cm}

\begin{minipage}[]{5.5cm}
Dans un repère orthonormé \Oijk {} on considère

\begin{itemize}
\item [$\bullet$] le point A de coordonnées (1~;~3~;~2),

 \item[$\bullet$]le vecteur $\vv{u}$ de coordonnées $\begin{pmatrix} 1\\1\\0\\\end{pmatrix}$

\item[$\bullet$] la droite $d$ passant par l'origine O du repère et admettant pour vecteur directeur $\vv{u}$. 
\end{itemize}




\end{minipage}
\begin{minipage}[]{4.5cm}
\psset{coorType=2,unit=1.5cm}
\begin{pspicture}(-2.5,-2.5)(4,3)
\pstThreeDCoor[xMin=0,xMax=3,yMin=-1.5,yMax=3.5,zMin=-1.5,zMax=2.5,IIIDticks]
\pstThreeDDot(1,3,2)\pstThreeDDot(1,3,0)\pstThreeDDot(2.5,2.5,0)\pstThreeDPut(2.5,2,0){$M_0$}
\pstThreeDPut(1,3.1,2.2){A}\pstThreeDPut(1,3,-0.3){A$'$}
\pstThreeDLine[linewidth=1.25pt]{->}(0,0,0)(1,1,0)\pstThreeDLine[linewidth=0.5pt]{-}(-2,-2,0)(4,4,0)
\pstThreeDLine[linewidth=1pt]{-}(2.5,2.5,0)(1,3,2)
\pstThreeDLine[linewidth=1pt]{->}(0,0,0)(1,3,2)
\pstThreeDLine[linewidth=1pt]{->}(0,0,0)(1,0,0)
\pstThreeDLine[linewidth=1pt]{-}(1,3,0)(1,3,2)
\pstThreeDLine[linewidth=1pt]{-}(2.5,2.5,0)(1,3,0)
\pstThreeDLine[linewidth=1pt,linestyle=dashed]{->}(0,0,0)(1,3,0)
\pstThreeDLine[linewidth=1pt]{->}(0,0,0)(0,0,1)
\pstThreeDLine[linewidth=1pt]{->}(0,0,0)(0,1,0)
\pstThreeDLine[linewidth=1pt,linestyle=dashed]{-}(0,0,0)(0,2.5,0)
\pstThreeDPut(1,-0.25,0){$\vv{\imath}$}\pstThreeDPut(0,1,0.15){$\vv{\jmath}$}
\pstThreeDPut(0,-0.2,1){$\vv{k}$}\pstThreeDPut(0,-0.1,-0.15){O}
\pstThreeDPut(0.9,0.65,0){$\vv{u}$}
\end{pspicture}

\end{minipage}


 \begin{enumerate}
\item 
$M(x~;~y~;~z) \in (d) \iff \vect{\text{O}M} = t\vect{u}, \, \text{avec } t \in \R$, soit :
$\left\{\begin{array}{l c l}
x&=&t\\
y&=&t\\
z&=&0
\end{array}\right., \, t \in \R$.

\item 
\begin{enumerate}
\item  
De $\vect{\text{A}M}\begin{pmatrix}t - 1\\t - 3\\0 - 2\end{pmatrix}$, on calcule :

A$M^2 = (t - 1)^2 + (t - 3)^2 + (- 2)^2 = t^2 + 1 - 2t + t^2 + 9 - 6t + 4 = 2t^2 - 8t+ 14.$


\item 


$2t^2 - 8t+ 14 = 2\left(t^2 - 4t + 7\right) = 2\left[(t - 2)^2 -  4 + 7 \right] = 2\left[(t - 2)^2 + 3\right]$.

La plus petite valeur de ce trinôme est obtenue quand le carré est nul, soit pour $t = 2$.

On a: $2t^2 - 8t+ 14 \geqslant 6$, soit $\text{A}M^2 \geqslant 6 \Rightarrow \text{A}M \geqslant \sqrt{6}$.

La plus petite distance est A$M_0 = \sqrt{6}$ avec $M_0(2~;~2~;~0)$.
\end{enumerate}
\item 
On a: $\vect{\text{A}M_0}\begin{pmatrix}1\\- 1\\0\end{pmatrix}$ et $\vect{u}\begin{pmatrix}1\\ 1\\0\end{pmatrix}$ est un vecteur directeur de $(d)$.

On a: $\vect{\text{A}M_0} \cdot \vect{u} = 1 - 1 + 0 = 0$ : les vecteurs sont orthogonaux donc les droites $\left(\text{A}M_0\right)$ et $d$ sont orthogonales.
\item 


$\vect{u}$ est orthogonal au plan horizontal d'équation $z = 0$. Comme A$'$ et $M_0$ appartiennent à ce plan le vecteur $\vect{u}$ est orthogonal au vecteur $\vect{\text{A}'M_0}$.

Donc le vecteur $\vect{u}$ est orthogonal à deux vecteurs non colinéaires du plan $\left(\text{AA}'M_0\right)$,donc la droite $(d)$ est orthogonale au plan 
$\left(\text{AA}'M_0\right)$. Le point $M_0$ est donc le projeté orthogonal de O sur le plan $\left(\text{AA}'M_0\right)$, donc O$M_0$ est la distance la plus courte du point O au plan $\left(\text{AA}'M_0\right)$.

\item 




Aire de la base AA$'M_0$ : on a AA$' = 2$ et A$'M_0^2 = (2 - 1)^2 + (2 - 3)^2 + 0^2 = 1 + 1 = 2$. D'où A$'M_0 = \sqrt{2}$.

On a donc $\mathcal{A}\left(\text{AA}'M_0\right) = \dfrac{2 \times \sqrt{2}}{2} = \sqrt{2}$.

D'autre part : O$M_0^2 = 2^2 + 2^2 = 8$, d'où O$M_0 = \sqrt{8} = \sqrt{4 \times 2} = 2\sqrt{2} = h$.

Finalement $V  = \dfrac{\sqrt{2} \times 2\sqrt{2}}{3} = \dfrac{4}{3}$.
\end{enumerate}

\vspace{0.75cm}

\end{document}
